\documentclass[../../main.tex]{subfiles}

\begin{document}

\section{Improving Large Language Models}

Up until now, we have discussed fundamental terms and very high-level
overview of pre-training phase. This section discusses very high-level
introduction to how a raw autocomplete model transitions to a model
that we use in daily life.

\subsection{Introduction to Post-Training}

One of the most important phase of training that makes the
chatbots chatbots are post-training. When a base model is
developed through pre-training phase, it is simply an autocomplete
model. Whereas, an aligned model is capable of generating
conversation-like responses. For instance, if the input is
``Why are dogs so cute?'', the base model will generate responses
like ``What are cats doing?'', which does not feel like a
conversation. Whereas, an aligned model will generate responses
like ``People find dogs cute because humans developed a biological
phenomena called baby schema.''

The post-training process contains two primary phases: Supervised
Fine-Tuning (SFT) and Reinforcement Learning from Human Feedback
(RLHF) On the high-level, SFT is the process where human provides
prompt and response samples. The model then learns from the samples
to provide a chatbot-like experience. After the phase, the model
can be reinforced through RLHF. There are multiple ways of doing
this, but a common method is making another ``reward model'' to
judge and provide feedback on a model's response. Instead of a
human manually judging and rewarding every responses, reward
model will predict how a human expert will reward (usually an
integer from 0 to 5) a response to more effectively align the model.
The specific process is quite out of scope for this note, and will
be introduced in the following notes.

Most of the ``learning'' happens in pre-training. However,
alignment is necessary to improve HHH, or helpfulness, honesty,
and harmlessness. And yes, SFT and RLHF are popular, but they
are not the only methods. Few other examples contain
\href{https://arxiv.org/abs/2305.11206}{LIMA},
\href{https://arxiv.org/abs/2212.08073}{RLAIF},
\href{https://arxiv.org/abs/2305.18290}{DPO},
\href{https://arxiv.org/abs/1707.06347}{PPO}, and
\href{https://arxiv.org/abs/2312.01552}{URIAL}.

\subsection{Reasoning}

As you probably have noticed, LLMs are predictions models.
In other words, they do not ``reason'' like how us humans
reason. Language models will predict the next likely token
that will provide a plausible response to a reasoning task.
One way to test if the models can actually reason is by
using benchmarks. One of the foundational datasets, which
is considered outdated now, is
\href{https://huggingface.co/datasets/openai/gsm8k}{GSM8K}
by OpenAI. This dataset is used to evaluate if a model can
reason grade school level math. The modern, highly challenging,
benchmark includes \href{https://epoch.ai/frontiermath}{FrontierMath}.
As you have noticed, benchmarking and reasoning is a giant
hurdle for many LLMs. Modern engineers classify reasoning
as another section in post-training to improve a model's logic.

\subsubsection{Prompting}

One of the most effective way to improve LLM reasoning is
to use prompting strategies.

\begin{definition}{}{}
One technique is \textbf{K-shot prompting} where you provide
$k$ examples for the model to follow a template. Zero-shot
prompting will be when $k = 0$, where you provide no examples,
but instructions. The second prompting strategy is \textbf{
\href{https://arxiv.org/abs/2201.11903}{Chain-of-Thought (CoT)
prompting}}. This strategy explicitly tells a model to break
a complex reasoning tasks into intermediate steps. 
\end{definition}

Let's take a look at an example on how K-shot prompting and
CoT prompting work together.

\begin{example}[Model Input]
\textbf{Instruction:}

Consider the following examples to provide a response to the
final equation. Provide a step by step explanation for your
answer as demonstrated in the samples.

\textbf{Example 1:}

Q: Henry the dog loves tomatoes. He originally had $5$ tomatoes.
Henry ate $3$ of them, and sneaked into the owner's yard to
secretly bring $4$ more. How many tomatoes does Henry have now?

A: Henry started with $5$, but ate $3$. Therefore, he has
$5 - 3 = 2$ tomatoes left. Because he brought 4 more from the
owner's yard, he now has $2 + 4 = 6$ tomatoes. Henry has $6$
tomatoes now.

\textbf{Example 2:}

Q: Bob the yak loves burgers. One day, his lizard friend Sofia
told him that she can bring three boxes of burgers than contain
2 burgers each. Bob instantly pleaded Sofia, and he received the
boxes. Out of hunger, he ate three burgers. How many burgers
does he have now?

A: It is assumed that Bob does not have any burgers before the
interaction with Sofia. With Sofia's virtue, Bob now has
$2 \cdot 3 = 6$ burgers. Because he ate three burgers from
six in total, Bob now has $6 - 3 = 3$ burgers left.

\textbf{Question:}

Charlotte the cat has five close friends. She wishes to split
her $120 \text{ kg}$ tuna equally with her friends. After the
allocation, Charlotte ate $4 \text{ kg}$ of tuna. How much
tuna in kg does Charlotte have left?
\end{example}

\begin{example}[Model Output]
Charlotte initially has $120 \text{ kg}$ of tuna. Because she
split it with her friends, she has $\frac{120}{5 + 1} = 20
\text{ kg}$ of tuna after the allocation. Since Charlotte ate
$4 \text{ kg}$ of tuna from $20 \text{ kg}$, she now has
$20 - 4 = 16 \text{ kg}$ of tuna.
\end{example}

As demonstrated, CoT has the potential to increase reasoning
of LLMs. However, it is important to note that although performance
may significantly ameliorate, numerous potential consequences
follow. For instance, because the model includes the intermediate
steps, the token generation considerably increases, which will
increase the cost. Moreover, when model starts to hallucinate
in the middle of the reasoning, the entire output from the
hallucination will no longer be useful. Furthermore, models
sometimes ``lie''. What I mean to lie is that some models tend
to provide a ``plausible CoT'' without actually providing one.
This is more related to AI safety, and let's save that for final
notes.

\subsubsection{Self-Consistency}

Another technique that could improve reasoning and performance
is \textbf{self-consistency}.

\begin{definition}{}{}
\textbf{Self-consistency} is a technique where a model generates
numerous possible outputs to an input, and aggregate the outputs
for the final output.
\end{definition}

The primary motivation of the technique is that just as mathematical
problems have numerous paths to an answer, we could derive an
``accurate'' output with different reasoning. However, like any
other techniques, there are limitations. One somewhat self-evident
caveat is that this requires much greater computation. Because the
technique requires multiple reasoning and aggregation for the final,
the amount of computation drastically increase. One other issue is
that this self-consistency only works on closed-solution prompts.
I mean different people and different outputs can have differing
opinions. Although this technique is not fully outdated, modern ML
engineers use more effective and efficient techniques.

The field of reinforcement learning is one of the most rapidly
developed field in ML. For instance, unlike RLHF, the emerging
method Reinforcement Learning with Verifiable Rewards (RLVR)
allows models to align themselves with minimum human contributions.
In the next note, we will discuss about neural network from scratch!

\end{document}


